\begin{figure}[ht]
    \centering
    \begin{subfigure}[t]{0.5\textwidth}
        \centering
        \begin{tikzpicture}[scale=1.5]
            \draw (0,0) circle (1.0);
            \draw (-2,-0.5) -- (2,0.5);
            \filldraw (-1.8,-0.45) circle (0.05);
            \draw[-{Latex[width=3]}] (-1.8,-0.45) -- (-1.4,-0.35);

            \filldraw[red] (-0.96,-0.24) circle (0.05);
            \filldraw[red] (0.96,0.24) circle (0.05);

            \node[] at (-1.8,-0.8) {A};
            \node[] at (-1.1,-0.64) [font=\small] {1};
            \node[] at (1.1,0.64) [font=\small] {2};
        \end{tikzpicture}
        \caption{}
    \end{subfigure}%
    ~
    \begin{subfigure}[t]{0.5\textwidth}
        \centering
        \begin{tikzpicture}[scale=1.5]
            \draw (0,0) circle (1.0);
            \draw (-2,-0.5) -- (2,0.5);
            \filldraw (-0.25,-0.07) circle (0.05);
            \draw[-{Latex[width=3]}] (-0.2,-0.05) -- (0.2,0.05);

            \filldraw[red] (-0.96,-0.24) circle (0.05);
            \filldraw[red] (0.96,0.24) circle (0.05);

            \node[] at (-0.25,-0.47) {B};
            \node[] at (-1.1,-0.64) [font=\small] {1};
            \node[] at (1.1,0.64) [font=\small] {2};
        \end{tikzpicture}
        \caption{}
    \end{subfigure}
    \caption{Intersections of a cylinder with a line from two different positions A and B shown in (a) and (b) respectively. The red points are the intersections labled with their stable indices 1 and 2.}
    \label{fig:dev/nev/cylinder-intersections}
\end{figure}
